% !TeX root = ../main.tex
\section*{Цель работы}
Изучение базовых принципов синтеза работы квадратурного приемника с цифровыми фильтрами дециматоров.

\section{Структурная схема устройства}

В качестве основной системы фильтрации (для пунктов 2--6) используется каскадное соединение фильтра Хогенауэра (CIC-дециматора) и КИХ-корректора. Схема включает в себя следующие этапы обработки:
\begin{itemize}
    \item Оцифровка входного сигнала с помощью аналого-цифрового преобразователя (АЦП) на частоте $f_s$;
    \item Разделение на синфазную (I) и квадратурную (Q) составляющие путем перемножения с сигналами цифрового гетеродина (ЦГ);
    \item Фильтрация каскадом из $N=8$ цифровых интеграторов;
    \item Понижение частоты дискретизации (децимация) в $R=10$ раз;
    \item Фильтрация каскадом из $N=8$ гребенчатых фильтров;
    \item Коррекция амплитудно-частотной характеристики в полосе пропускания с помощью КИХ-фильтра компенсатора (ФК).
\end{itemize}

\begin{figure}[H]
    \centering
    \includegraphics[width=0.9\textwidth]{CIC + КИХ-корректор.PNG}
    \caption{Структурная схема устройства цифровой фильтрации (CIC + КИХ-корректор)}
    \label{fig:struct_main}
\end{figure}

\begin{figure}[H]
    \centering
    \begin{minipage}[t]{0.48\textwidth}
        \vspace{0pt}
        \centering
        \includegraphics[width=\linewidth]{int.PNG}
        \caption{Цифровой интегратор}
     \end{minipage}
    \hfill
    \begin{minipage}[t]{0.48\textwidth}
        \vspace{0pt}
        \centering
        \includegraphics[width=\linewidth]{greb.PNG}
        \caption{Цифровой гребенчатый фильтр}
    \end{minipage}
\end{figure}

Для выполнения пункта 7 каскад интеграторов, гребенчатых фильтров и КИХ-корректор заменяются на единый КИХ-фильтр дециматор. В этом случае понижение частоты дискретизации в 10 раз осуществляется непосредственно на выходе КИХ-фильтра.

\begin{figure}[H]
    \centering
    \includegraphics[width=0.9\textwidth]{FIR.PNG}
    \caption{Структурная схема устройства цифровой фильтрации (единый КИХ-дециматор)}
    \label{fig:struct_fir}
\end{figure}
